\documentclass[11pt,a4paper]{article}
\usepackage[T1]{fontenc}
\usepackage[utf8]{inputenc}
\usepackage[ngerman]{babel}
\usepackage[margin=2.5cm]{geometry}
\usepackage[hidelinks]{hyperref}
\usepackage{array}
\setlength{\parindent}{0pt}
\setlength{\emergencystretch}{2em}
\renewcommand{\arraystretch}{1.25}
\sloppy
\begin{document}
\shorthandoff{"}
\begin{center}
  {\Large\textbf{Webtechnologien -- Übungsblatt 6}}
\end{center}
\section*{Aufgabe 2: Simple Node.js Product List API}
\textbf{Aufgabe:} Produkte aus \texttt{products.json} laden und über
Express als JSON anbieten: einmal alle IDs, einmal ein vollständiges
Produkt. Eine unbekannte ID soll den Status 404 liefern.

\medskip
\textbf{Ergebnis:} In \texttt{server.js} sind die Platzhalter ergänzt.
Die unangetastete Vorlage bleibt als \texttt{server\_ORIGINAL.js} erhalten;
\texttt{products.json} bleibt unverändert.
\begin{itemize}
  \item \texttt{fs.readFileSync(\_\_dirname + '/products.json', 'utf8')}
    liest die Datei einmal beim Start als Text. \texttt{\_\_dirname} ist
    der Ordner der JavaScript-Datei; der Dateipfad hängt dadurch nicht
    vom aktuellen Arbeitsordner ab.
  \item \texttt{JSON.parse} macht daraus ein Array von Produktobjekten.
    \texttt{console.log(products)} zeigt es zur Kontrolle im Terminal.
    Das synchrone Lesen ist hier einfach, weil es nur beim Start erfolgt.
  \item Die Startseite erhält wie in Aufgabe 1 den Inhaltstyp HTML.
    \texttt{/products} nimmt mit \texttt{map} aus jedem Produkt nur die ID.
    Die Reihenfolge der Datei bleibt erhalten. \texttt{res.json} sendet
    das Objekt mit \texttt{productIds} und setzt den JSON-Inhaltstyp.
  \item \texttt{/product/:id} liest \texttt{req.params.id} und sucht
    mit \texttt{find} das passende Objekt. Bei Erfolg kommt das gesamte
    Objekt mit \texttt{id}, \texttt{title}, \texttt{desc} und \texttt{price}.
    Ohne Treffer beendet \texttt{return} die Route mit 404 und einer
    JSON-Fehlermeldung; es wird keine zweite Antwort gesendet.
\end{itemize}

\textbf{Starten:} Im Ordner \texttt{aufgabe2} zuerst einmal
\texttt{npm install} ausführen.
Mit \texttt{node server.js} starten und
\url{http://localhost:3000/products} öffnen. Ein einzelnes Produkt ist unter
\url{http://localhost:3000/product/sony24105} erreichbar.
Vorher andere Übungsserver auf Port 3000 beenden; Strg+C stoppt den Server.

\medskip
\textbf{Test:} Mit Node.js v24 gestartet, Anfragen mit \texttt{curl}; alle Fälle liefern genau diese Antworten (JSON):
\begin{center}
\begin{tabular}{|>{\raggedright\arraybackslash}p{5.2cm}|>{\raggedright\arraybackslash}p{8.5cm}|}
\hline
\textbf{Fall} & \textbf{Ergebnis} \\ \hline
GET \texttt{/products} & 200; \texttt{\{"productIds":["sonyAIV", "sony24105", "NiconDX16400", "tamron1724"]\}} \\ \hline
GET \texttt{/product/sony24105} & 200; \texttt{\{"id":"sony24105", "title":"Sony 24-105mm lens", "desc":"Best Sony travel zoom", "price":1300\}} \\ \hline
GET \texttt{/product/nope} & 404; \texttt{\{"error":"product not found"\}} \\ \hline
\end{tabular}
\end{center}

\medskip
{\small \textbf{Hilfsmittel:} Claude Code (Claude Opus 5.5) und OpenAI Codex (GPT-6.1-Sol) für Erstellung, Prüfung und Dokumentation.}
\end{document}
